\documentclass[11pt,a4paper]{article}
\usepackage[utf8]{inputenc}
\usepackage{amsmath, amssymb, amsfonts}
\usepackage{booktabs}
\usepackage{geometry}
\usepackage{hyperref}

\geometry{margin=1in}

\title{\textbf{A Geometric and Bayesian Formulation of Autoregressive Language Models}}
\author{\textbf{Akshay Mishra}}
\date{\today}

\begin{document}
	
	\maketitle
	
	\begin{abstract}
		This paper formalizes the mechanics of autoregressive Large Language Models (LLMs) by reframing their forward pass and pre-training objectives through the lens of Bayesian inference, high-dimensional geometry, and statistical mechanics. We show that the conditional probability distribution generated by the Transformer architecture can be explicitly decomposed into an intrinsic token prior and a contextual likelihood. We map these probabilistic components directly onto the network's structural blocks—specifically Multi-Head Attention (MHA) and Feed-Forward Networks (FFN)—and demonstrate how the final output token distribution is modulated by the statistical variance of the latent residual stream, acting as a dynamic inverse temperature. Finally, we formalize pre-training as a self-organizing geometric optimization that resolves these components into a stable semantic coordinate system.
	\end{abstract}
	
	
	
	\section{Introduction and the Bayesian Framing}
	
	An autoregressive large language model operates over a finite vocabulary $V$~\cite{bengio2003neural}. Given a context sequence of $n$ tokens, $C = (t_1, t_2, \dots, t_n)$, where $t_i \in V$, the model defines a discrete conditional probability distribution over the vocabulary for the next token $w \in V$:
	
	\begin{equation}
		P(w \mid C)
	\end{equation}
	
	While modern deep learning architectures evaluate this distribution directly using discriminative causal language modeling~\cite{Radford2019LanguageMA}, we can invert this formulation using Bayes' theorem to expose the underlying generative dynamics~\cite{jaynes1957}:
	
	\begin{equation}
		P(w \mid C) = \frac{P(C \mid w) \cdot P(w)}{P(C)}
	\end{equation}
	
	Because the denominator $P(C) = \sum_{v \in V} P(C \mid v)P(v)$ represents the marginal likelihood of the context sequence, it acts as a constant normalization factor across all candidate tokens at step $n+1$. Thus, the relative probability mass assigned to any token $w$ is determined strictly by the numerator:
	
	\begin{equation}
		P(w \mid C) \propto P(C \mid w) \cdot P(w)
	\end{equation}
	
	Where:
	\begin{itemize}
		\item $P(w)$ is the \textbf{unconditional prior probability} of the token $w$ occurring in the language, independent of local context.
		\item $P(C \mid w)$ is the \textbf{contextual likelihood}, measuring how strongly the historical sequence $C$ retroactively aligns with or necessitates the trailing token $w$.
	\end{itemize}
	
	\section{Mathematical Decomposition of Likelihood \texorpdfstring{$P(C \mid w)$}{P(C | w)}}
	
	In a Transformer architecture~\cite{vaswani2017attention}, the context sequence $C$ is transformed into a sequence of continuous vectors. Let $X \in \mathbb{R}^{n \times d}$ represent the input embedding matrix for the context, where $d$ is the model dimension. The final hidden state at the terminal position $n$ after $L$ layers is denoted as $h_n^{(L)} \in \mathbb{R}^d$. 
	
	The unnormalized log-probabilities (logits) $z_w$ for each token $w$ are computed by projecting $h_n^{(L)}$ onto an unembedding matrix $U \in \mathbb{R}^{|V| \times d}$, such that $z_w = h_n^{(L)} \cdot u_w$, where $u_w \in \mathbb{R}^d$ is the target token's unembedding vector~\cite{press2017using}. 
	
	The log-likelihood $\log P(C \mid w)$ is shaped by the interaction of two foundational architectural components within each layer $l$: Multi-Head Attention and the position-wise Feed-Forward Network.
	
	\subsection{Multi-Head Attention as Contextual Routing}
	The Multi-Head Attention block at layer $l$ and position $n$ computes a dynamically weighted combination of all preceding token states~\cite{vaswani2017attention}. Let $H$ be the number of attention heads, and $d_k$ be the head dimension. For a single attention head $g \in \{1, \dots, H\}$, let $W_Q^{(l, g)}, W_K^{(l, g)}, W_V^{(l, g)} \in \mathbb{R}^{d \times d_k}$ be the Query, Key, and Value projection matrices, respectively. 
	
	The terminal query vector for head $g$ is defined as $q_n^{(l, g)} = h_n^{(l-1)}W_Q^{(l, g)}$, while the keys and values for all historical positions $i \in \{1, \dots, n\}$ are $k_i^{(l, g)} = h_i^{(l-1)}W_K^{(l, g)}$ and $v_i^{(l, g)} = h_i^{(l-1)}W_V^{(l, g)}$. The dynamic attention weights $\alpha_{n,i}^{(l, g)}$ formalize the conditional token-to-token dependencies:
	
	\begin{equation}
		\alpha_{n,i}^{(l, g)} = \frac{\exp\left(\frac{q_n^{(l, g)} \cdot k_i^{(l, g)}}{\sqrt{d_k}}\right)}{\sum_{j=1}^n \exp\left(\frac{q_n^{(l, g)} \cdot k_j^{(l, g)}}{\sqrt{d_k}}\right)}
	\end{equation}
	
	The concatenated outputs of all heads are projected back into the residual stream via the global output matrix $W_O^{(l)} \in \mathbb{R}^{d \times d}$. By partitioning $W_O^{(l)} = [ (W_{O, 1}^{(l)})^T \dots (W_{O, H}^{(l)})^T ]^T$, where each head's output projection slice is $W_{O, g}^{(l)} \in \mathbb{R}^{d_k \times d}$, the total attention contribution to the log-likelihood of a target token $w$ decomposes into a sum of bilinear forms:
	
	\begin{equation}
		\log P(C \mid w)_{\text{MHA}} \sim \sum_{g=1}^H \sum_{i=1}^n \alpha_{n,i}^{(l, g)} \left( h_i^{(l-1)} W_V^{(l, g)} W_{O, g}^{(l)} u_w \right)
	\end{equation}
	
	This isolates two distinct mechanisms: $\alpha_{n,i}^{(l, g)}$ handles the contextual routing (determining which past tokens matter), while the dot product acts as a semantic alignment function between the historical features of token $i$ and the target token $w$.
	
	\subsection{Feed-Forward Networks as Concept Memory Retrieval}
	Following the attention block, the updated hidden state $x_n^{(l)} = h_n^{(l-1)} + a_n^{(l)}$ enters the position-wise Feed-Forward Network (FFN). The FFN operates as a key-value associative memory~\cite{geva-etal-2021-transformer}. Let $W_1^{(l)} \in \mathbb{R}^{d \times d_{ff}}$ represent a bank of intermediate keys $\{k_m^{(l)}\}_{m=1}^{d_{ff}}$, and $W_2^{(l)} \in \mathbb{R}^{d_{ff} \times d}$ represent a bank of intermediate values $\{v_m^{(l)}\}_{m=1}^{d_{ff}}$.
	
	Applying a non-linear activation function $\sigma$ (e.g., GELU~\cite{hendrycks2016gelu}), the operation at the terminal position is:
	
	\begin{equation}
		\text{FFN}(x_n^{(l)}) = \sum_{m=1}^{d_{ff}} \sigma\left( x_n^{(l)} \cdot k_m^{(l)} + b_{1,m}^{(l)} \right) v_m^{(l)}
	\end{equation}
	
	Projecting this operation onto the target token embedding $u_w$ reveals the FFN's mathematical formulation of the likelihood:
	
	\begin{equation}
		\log P(C \mid w)_{\text{FFN}} \sim \sum_{m=1}^{d_{ff}} \sigma\left( x_n^{(l)} \cdot k_m^{(l)} + b_{1,m}^{(l)} \right) \cdot (v_m^{(l)} \cdot u_w)
	\end{equation}
	
	Here, $\sigma\left( x_n^{(l)} \cdot k_m^{(l)} + b_{1,m}^{(l)} \right)$ scales how strongly the synthesized context satisfies the $m$-th latent conceptual rule, while $(v_m^{(l)} \cdot u_w)$ dictates the historical co-occurrence correlation between that rule and the literal vocabulary token $w$.
	
	\section{Geometric Formulation of the Prior \texorpdfstring{$P(w)$}{P(w)} and Logit Mechanics}
	
	The exact computation of the unconditional prior $P(w)$ requires marginalizing over the infinite space of all possible context sequences $\mathcal{C}$:
	
	\begin{equation}
		P(w) = \sum_{C \in \mathcal{C}} P(w \mid C)P(C)
	\end{equation}
	
	Transformers approximate this expectation statically within their parameters via two geometric mechanics:
	
	\subsection{Unembedding Bias \texorpdfstring{($b_U$)}{(b\_U)}}
	In architectures featuring an explicit unembedding bias $b_U \in \mathbb{R}^{|V|}$, the logit calculation is defined as $z_w = (h_n^{(L)} \cdot u_w) + b_{U,w}$. Under an uninformative context where $h_n \to \mathbf{0}$, the logit collapses directly to the scalar bias:
	
	\begin{equation}
		P(w) \approx \frac{\exp(b_{U,w})}{\sum_{v \in V} \exp(b_{U,v})}
	\end{equation}
	
	\subsection{Embedding Norms \texorpdfstring{($\|u_w\|$)}{(||u\_w||)}}
	In bias-free architectures, the prior is encoded in the Euclidean magnitude of the token vector, $\|u_w\|$. Decomposing the final dot product geometrically yields:
	
	\begin{equation}
		z_w = \|h_n^{(L)}\| \|u_w\| \cos(\theta_w)
	\end{equation}
	
	Where $\theta_w$ is the high-dimensional angle between the final context state and the token vector. Tokens with high baseline frequencies in the training corpus develop longer vector magnitudes $\|u_w\|$ during optimization, amplifying their baseline logit regardless of the contextual direction $\cos(\theta_w)$~\cite{gao2019representation}.
	
	\section{Statistical Mechanics of the Latent Stream}
	
	The remaining term in the geometric logit decomposition is the norm of the final hidden state, $\|h_n^{(L)}\|$. This component functions as a dynamic, context-dependent entropy controller.
	
	\subsection{Activation Variance as Inverse Temperature}
	Let $h_n = [x_1, x_2, \dots, x_d] \in \mathbb{R}^d$ be a row vector representing the final hidden state $h_n^{(L)}$. If we treat the individual channel activations as samples from an internal distribution with a mean driven to zero by layer normalization ($\mu_x \approx 0$)~\cite{ba2016layernorm}, the cross-channel variance is:
	
	\begin{equation}
		\sigma^2 = \frac{1}{d} \sum_{i=1}^d (x_i - \mu_x)^2 \approx \frac{1}{d} \sum_{i=1}^d x_i^2
	\end{equation}
	
	Because the squared Euclidean norm is $\|h_n\|^2 = \sum_{i=1}^d x_i^2$, the vector magnitude can be expressed as a function of activation variance:
	
	\begin{equation}
		\|h_n\| = \sqrt{d \cdot \sigma^2}
	\end{equation}
	
	When substituting this back into the vocabulary-wide Softmax function, $\|h_n\|$ maps perfectly to the \textbf{inverse temperature} parameter $\beta = \frac{1}{T}$ from statistical mechanics~\cite{jaynes1957, ackley1985learning}:
	
	\begin{equation}
		P(w \mid C) = \frac{\exp\big(\beta \cdot \|u_w\| \cos(\theta_w)\big)}{\sum_{v \in V} \exp\big(\beta \cdot \|u_v\| \cos(\theta_v)\big)}, \quad \text{where } \beta = \|h_n\|
	\end{equation}
	
	\begin{itemize}
		\item \textbf{High Contextual Certainty:} When the context contains a highly predictable pattern, successive attention heads and FFN layers write aligned information into the residual stream. This constructive interference inflates the cross-channel variance $\sigma^2$, causing $\|h_n\|$ to expand. This yields a high $\beta$, sharpening the distribution and lowering its entropy.
		\item \textbf{High Contextual Ambiguity:} When the context is entropic or transitions between distinct topics, architectural sub-layers write orthogonal or destructive vectors into the stream. The variance $\sigma^2$ collapses, shrinking $\|h_n\|$. This yields a low $\beta$, flattening the distribution into a high-entropy state.
	\end{itemize}
	
	\subsection{Cosine Similarity as a Correlation Coefficient}
	Similarly, treating the dimensions of $h_n$ and $u_w$ as zero-mean random variables allows us to map the likelihood component $\cos(\theta_w)$ directly to the \textbf{Pearson correlation coefficient ($r$)} in feature space:
	
	\begin{equation}
		\cos(\theta_w) = \frac{\text{Cov}(h_n, u_w)}{\sigma_{h_n} \sigma_{u_w}} = r_{h_n, u_w}
	\end{equation}
	
	where $\sigma_{h_n}$ and $\sigma_{u_w}$ denote the standard deviations of the components of $h_n$ and $u_w$, respectively. Thus, the unnormalized logit is exactly proportional to the total feature co-variance scaled by the model dimension: $z_w = \text{Cov}(h_n, u_w) \cdot d$.
	
	\section{Statistical Description of Pre-Training}
	
	During inference, $u_w$ is static. During pre-training, both the context state $h_n$ and the vocabulary matrix $U$ are discovered simultaneously via the minimization of the cross-entropy loss~\cite{shannon1948mathematical} for a target next-token $w^*$:
	
	\begin{equation}
		\mathcal{L} = -\log \frac{\exp(h_n \cdot u_{w^*})}{\sum_{v \in V} \exp(h_n \cdot u_v)}
	\end{equation}
	
	This framework can be understood as a dual-force geometric optimization process:
	
	\subsection{The Mutual Attraction Force}
	The gradient with respect to the correct token vector $u_{w^*}$ acts as an attractive force:
	
	\begin{equation}
		\frac{\partial \mathcal{L}}{\partial u_{w^*}} = -\big(1 - P(w^* \mid C)\big) h_n
	\end{equation}
	
	This force dynamically rotates $u_{w^*}$ and the context representation $h_n$ toward one another, maximizing their directional correlation $\cos(\theta_{w^*})$. Simultaneously, if a token $w^*$ is consistently validated across diverse contexts, its vector norm $\|u_{w^*}\|$ is driven upward to minimize loss across the corpus, mapping out its global prior $P(w)$.
	
	\subsection{The Global Repulsion Force}
	Conversely, the denominator of the Softmax generates a contrastive repulsion force against all incorrect tokens $v \neq w^*$, resembling contrastive representation learning frameworks~\cite{oord2018representation}:
	
	\begin{equation}
		\frac{\partial \mathcal{L}}{\partial u_v} = P(v \mid C) h_n
	\end{equation}
	
	This pushes the context state $h_n$ away from incorrect vocabulary vectors, forcing them toward orthogonality ($\cos(\theta_v) \to 0$).
	
	\subsection{Emergence of Distributional Semantics}
	Because the network minimizes this loss across billions of tokens, semantic clustering emerges naturally. If two distinct tokens $w_1$ and $w_2$ share high contextual substitutionality, they will repeatedly receive identical directional updates from shared context states $h_n$. This self-organizing system settles into a geometric steady-state where vocabulary vectors discover a coordinate system whose spatial proximity mirrors latent semantic relationships, validating the distributional hypothesis~\cite{harris_distributional_1954}:
	
	\begin{equation}
		\cos(\theta_{w_1, w_2}) \to 1 \iff \mathcal{C}_{w_1} \approx \mathcal{C}_{w_2}
	\end{equation}
	
	where $\mathcal{C}_w$ represents the distribution of contexts in which token $w$ appears.
	
	\section{Conclusion}
	
	By formalizing the Large Language Model as a Bayesian engine, we transition from a qualitative understanding of neural network behaviors to an exact geometric framework. Multi-Head Attention and Feed-Forward Networks act in tandem to iteratively construct a context vector whose final magnitude represents the model's statistical certainty ($\beta$), and whose orientation represents structural feature configurations. The final token prediction is thus a direct measure of the Pearson correlation between this dynamic contextual variance and the self-organized coordinate landscape of the vocabulary embeddings.
	
	\bibliographystyle{plain}
	\begin{thebibliography}{99}
		
		\bibitem{vaswani2017attention}
		Vaswani, A., Shazeer, N., Parmar, N., Uszkoreit, J., Jones, L., Gomez, A. N., Kaiser, \L{}. and Polosukhin, I. (2017).
		Attention is all you need.
		\textit{Advances in neural information processing systems}, 30.
		
		\bibitem{bengio2003neural}
		Bengio, Y., Ducharme, R., Vincent, P. and Janvin, C. (2003).
		A neural probabilistic language model.
		\textit{Journal of Machine Learning Research}, 3, pp.1137--1155.
		
		\bibitem{Radford2019LanguageMA}
		Radford, A., Wu, J., Child, R., Luan, D., Amodei, D. and Sutskever, I. (2019).
		Language models are unsupervised multitask learners.
		\textit{OpenAI blog}, 1(9), p.9.
		
		\bibitem{press2017using}
		Press, O. and Wolf, L. (2017).
		Using the output embedding to improve language models.
		\textit{Proceedings of the 15th Conference of the European Chapter of the Association for Computational Linguistics: Volume 2, Short Papers}, pp.157--163.
		
		\bibitem{geva-etal-2021-transformer}
		Geva, M., Schuster, R., Berant, J. and Levy, O. (2021).
		Transformer feed-forward layers are key-value memories.
		\textit{Proceedings of the 2021 Conference on Empirical Methods in Natural Language Processing}, pp.5484--5495.
		
		\bibitem{hendrycks2016gelu}
		Hendrycks, D. and Gimpel, K. (2016).
		Gaussian error linear units (GELUs).
		\textit{arXiv preprint arXiv:1606.08415}.
		
		\bibitem{gao2019representation}
		Gao, J., He, D., Tan, X., Qin, T., Wang, L. and Liu, T.Y. (2019).
		Representation degeneracy problem in training natural language generation models.
		\textit{International Conference on Learning Representations (ICLR)}.
		
		\bibitem{ba2016layernorm}
		Ba, J.L., Kiros, J.R. and Hinton, G.E. (2016).
		Layer normalization.
		\textit{arXiv preprint arXiv:1607.06450}.
		
		\bibitem{jaynes1957}
		Jaynes, E.T. (1957).
		Information theory and statistical mechanics.
		\textit{Physical review}, 106(4), p.620.
		
		\bibitem{ackley1985learning}
		Ackley, D.H., Hinton, G.E. and Sejnowski, T.J. (1985).
		A learning algorithm for Boltzmann machines.
		\textit{Cognitive Science}, 9(1), pp.147--169.
		
		\bibitem{shannon1948mathematical}
		Shannon, C.E. (1948).
		A mathematical theory of communication.
		\textit{The Bell System Technical Journal}, 27(3), pp.379--423.
		
		\bibitem{oord2018representation}
		van den Oord, A., Li, Y. and Vinyals, O. (2018).
		Representation learning with contrastive predictive coding.
		\textit{arXiv preprint arXiv:1807.03748}.
		
		\bibitem{harris_distributional_1954}
		Harris, Z.S. (1954).
		Distributional structure.
		\textit{Word}, 10(2-3), pp.146--162.
		
	\end{thebibliography}

\end{document}